\section{Experimental Setup}
\label{sec:setup}
The experimental testbed consists of three AI execution clients, a coordinating
server, and a passive network proxy positioned along the communication path.
The testbed is heterogeneous in both hardware and operating system: the clients
are an NVIDIA Jetson Orin Nano and a Dell desktop running Ubuntu, and a Dell
laptop running Windows~11, whereas the server, hosted on an NVIDIA Jetson AGX
Orin, and the proxy run Ubuntu. The clients therefore differ in compute
capability as well as in their network stacks; for example, the Windows client
omits the TCP timestamp option that the Ubuntu hosts use by default. In two
experiments, the Jetson AGX Orin also served as the embedded client. Clients
connect to the proxy, which relays their traffic to the server, so every
client--server exchange traverses the observation point. The proxy captures packet-level metadata, including packet
sizes, directions, timing, traffic volumes, and burst characteristics, without
accessing model parameters or application payloads. These quantities are
observable whether or not the payload is encrypted. Table~\ref{tab:ai_workloads}
summarizes the AI workloads, architectural configurations, and training
datasets. Each experiment executes 100 federated learning rounds. Federated
training uses a custom synchronous FedAvg implementation in PyTorch, in which
the server and clients exchange length-prefixed messages over TCP; in each
round, every client downloads the global model, trains locally, and uploads one
update before the server aggregates and advances to the next round. The
experimental corpus comprises 32 independent experiments and 96 client-level
aggregated observations. Traffic is represented using 0.5, 1, 2, and 5~s
observation windows and complete-trace aggregates to examine fingerprinting
performance across temporal scales. Classification is evaluated across the
family, architecture, and variant levels of the model hierarchy, with the
validation protocol detailed in Sec.~\ref{subsec:classification}.

\begin{table}[!htbp]
\centering
\caption{AI Workloads, Architectures, Variants, and Datasets}
\label{tab:ai_workloads}
\scriptsize
\renewcommand{\arraystretch}{1.10}
\setlength{\tabcolsep}{2pt}
\begin{tabular}{
>{\RaggedRight\arraybackslash}p{0.19\columnwidth}
>{\RaggedRight\arraybackslash}p{0.22\columnwidth}
>{\RaggedRight\arraybackslash}p{0.32\columnwidth}
>{\RaggedRight\arraybackslash}p{0.16\columnwidth}
}
\hline
\textbf{Family} &
\textbf{Architecture} &
\textbf{Representative Variants} &
\textbf{Dataset} \\
\hline
Autoencoder &
Dense, Conv., VAE &
Layer configurations, $\beta$-VAE, Conv. VAE &
Fashion-MNIST \\
CNN &
ResNet, MobileNet, EfficientNet &
ResNet18/34/50/101, MobileNetV2/V3,
EfficientNet-B0/B1/B2 &
Fashion-MNIST \\
RNN &
LSTM, GRU &
Single-layer, multilayer,
bidirectional variants &
UCI HAR \\
Transformer &
BERT, DistilBERT, Tiny Transformer &
BERT model sizes,
lightweight configurations &
AG News \\
\hline
\end{tabular}
\end{table}

\section{Methodology}
\label{sec:methodology}
The proposed hierarchical AI fingerprinting framework identifies AI model
characteristics from network side-channel observations without requiring access
to application payloads or model parameters. As illustrated in
Fig.~\ref{fig:architecture}, the framework consists of passive traffic
observation, multiscale feature extraction, hierarchical classification, and
fingerprint evaluation.

\subsection{Fingerprinting Dataset and Ground Truth}
\label{subsec:dataset}
The dataset comprises about 499 million packets (670~GB, 190~hours of traffic)
from 32 experiments, one per variant, each observed from three clients
(Table~\ref{tab:fp_dataset}). Each of the 96 client traces yields one
complete-trace observation and a sequence of windows; round-level quantities
are available for the 90 traces whose rounds are fully reconstructed. Because
larger models produce longer traces, training weights every experiment equally.
Each observation is described by 68 size-normalized features covering direction
and volume ratios, transmission rates, and packet-size, interarrival, burst, and
idle statistics. Complete traces add nine compute-timing features derived from
the round structure, giving 77 features per client and 231 in the cross-device
representation~\eqref{eq:crossdevice}. Labels and per-model metadata (parameter
count, serialized size, and training FLOPs per sample) are computed from the
model source code and are used only for evaluation.

\begin{table}[!htbp]
\centering
\caption{Fingerprinting Dataset Composition}
\label{tab:fp_dataset}
\scriptsize
\renewcommand{\arraystretch}{1.05}
\setlength{\tabcolsep}{4pt}
\begin{tabular}{lccccc}
\hline
\textbf{Family} & \textbf{Arch.} & \textbf{Var.} & \textbf{Traces} &
\textbf{5~s win.} & \textbf{0.5~s win.} \\
\hline
Autoencoder & 3  & 9  & 27 & 7{,}943  & 57{,}786 \\
CNN         & 3  & 9  & 27 & 35{,}947 & 298{,}965 \\
RNN         & 2  & 6  & 18 & 2{,}578  & 12{,}036 \\
Transformer & 3  & 8  & 24 & 39{,}585 & 345{,}817 \\
\hline
\textbf{Total} & \textbf{11} & \textbf{32} & \textbf{96} & \textbf{86{,}053} & \textbf{714{,}604} \\
\hline
\end{tabular}
\end{table}

\subsection{Data Preprocessing}
\label{subsec:preprocessing}
Raw captures are converted into analysis-ready traces in five steps.

\emph{Trace assembly.} For every client connection, the proxy records a sequence
of packet entries (timestamp, frame length, direction, and TCP flags). Long
experiments are captured in consecutive chunks, which are placed on a common
timeline using their recorded capture-start offsets, and each trace is mapped to
its client device through the client address logged by the proxy. Connection
attempts that complete no training round are discarded.

\emph{Payload recovery.} The TCP payload of each packet is obtained by
subtracting the Ethernet, IPv4, and TCP header length from the frame length.
The header length is inferred per connection from the size of pure
acknowledgments: 66 bytes when the TCP timestamp option is used and 54 bytes
when it is not, as on the Windows client.
Packets without payload, such as pure acknowledgments, are excluded from
transfer reconstruction.

\emph{Transfer and round reconstruction.} Transfers are recovered as maximal
runs of same-direction data packets, tolerating small opposite-direction control
packets (at most 300~B) and transient stalls caused by wireless retransmission,
which must not split a transfer. Because rounds alternate strictly between
download and upload, the model-carrying transfers are identified as the
byte-weighted dominant size class in each direction, and each upload is paired
with the preceding download to delimit a round.

\emph{Windowing and filtering.} Each trace is segmented into non-overlapping
windows on the global timeline (Sec.~\ref{subsec:multiscale}). Windows with
fewer than 20 packets are discarded, and near-idle windows are additionally
removed for cross-device window classification.

\emph{Normalization.} Features that scale directly with model size or trace
length (absolute byte and packet counts, durations, and connection and TCP flag
counts) are excluded, so that the remaining features describe the shape of the
traffic rather than its volume. Missing and non-finite values, which arise for
windows without traffic in one direction or for an absent device in the
cross-device representation, are set to zero for windows and to the feature
median for traces. Features are standardized using statistics computed on the
training folds only.

\subsection{Traffic Observation and Feature Extraction}
Network traffic exchanged between AI clients and the coordinating server is
observed at an intermediate proxy. Each captured packet is represented by
observable metadata, including its timestamp, size, and communication direction.
Let
\begin{equation}
p_i = (t_i, s_i, d_i)
\label{eq:packet}
\end{equation}
denote the $i$th observed packet, where $t_i$ is its arrival timestamp, $s_i$ is
its size, and $d_i$ indicates its communication direction. From the resulting
packet sequence, the framework extracts statistical features describing packet
sizes, interarrival times, directional traffic volumes, transmission rates, and
burst behavior.

Beyond per-packet statistics, we exploit the round structure of federated
training. Each round consists of a model download, a local-training interval
with no traffic, and a model upload. From the directional packet sequence we
recover, for round $k$, the end of the download $t^{\mathrm{d}}_{k}$ and the
start of the upload $t^{\mathrm{u}}_{k}$, and define the \emph{compute gap}
\begin{equation}
g_k = t^{\mathrm{u}}_{k} - t^{\mathrm{d}}_{k},
\label{eq:gap}
\end{equation}
a passive proxy for local training time that excludes the inter-client
synchronization wait. Together, these measurements form a network fingerprint
representing the communication and computation characteristics of the executing
AI workload. The per-trace summary augments the packet-morphology features with
the compute gap normalized by uploaded payload, its dispersion, and cross-device
timing ratios (Sec.~\ref{subsec:crossdevice}).

\subsection{Multiscale Traffic Representation}
\label{subsec:multiscale}
To investigate the temporal availability of AI fingerprints, the observed
traffic is represented at multiple time scales. For an observation duration
$\Delta$, packets are grouped into temporal windows, and a feature vector is
computed for each window:
\begin{equation}
\mathbf{x}_{k}^{(\Delta)}
=
\phi\left(
\{p_i \mid k\Delta \leq t_i < (k+1)\Delta\}
\right),
\label{eq:window_features}
\end{equation}
where $\phi(\cdot)$ denotes the traffic feature extraction function. The
framework considers observation durations $\Delta \in \{0.5,1,2,5\}$~s. In
addition, complete client traces are summarized into aggregated fingerprints to
characterize communication behavior over the entire observed execution.

A single short window observes only a fragment of one round. We therefore
aggregate a client's window-level posteriors to a trace-level decision by
averaging,
\begin{equation}
P(y \mid \mathcal{T})
= \frac{1}{|\mathcal{T}|}\sum_{k \in \mathcal{T}}
P\!\left(y \mid \mathbf{x}_k^{(\Delta)}\right),
\label{eq:pooling}
\end{equation}
where $\mathcal{T}$ indexes the windows of a trace. This recovers the
round-level signal lost to short windows without additional training, and
enables comparison between short-duration traffic signatures and complete-trace
observations.

For cross-device window classification, the three clients' windows covering
the same global interval are concatenated, analogously to~\eqref{eq:crossdevice}.
Near-idle windows, whose transferred bytes fall below $5\%$ of the trace's
busy-window volume or whose retransmission share exceeds $5\%$, are removed
before classification; this filter is fixed in advance and uses no labels.
For short windows, variant identity is assessed by nearest-reference matching on
the size of a model download that completes within the window.
To study observation times beyond the window length, the posteriors of
consecutive 5~s windows are averaged with~\eqref{eq:pooling} over
non-overlapping spans of 10 to 30~s.

\subsection{Hierarchical AI Model Fingerprinting}
Rather than treating all AI models as unrelated classes, the proposed framework
represents model identity through a three-level hierarchy:
\begin{equation}
\mathcal{H} =
\mathcal{F}
\rightarrow
\mathcal{A}
\rightarrow
\mathcal{V},
\label{eq:hierarchy}
\end{equation}
where $\mathcal{F}$, $\mathcal{A}$, and $\mathcal{V}$ denote model families,
architectures, and variants, respectively. Given a network fingerprint
$\mathbf{x}$, family identification is performed as
\begin{equation}
\hat{f}
=
\arg\max_{f \in \mathcal{F}}
P(f \mid \mathbf{x}).
\label{eq:family}
\end{equation}
Architecture identification is then conditioned on the selected family:
\begin{equation}
\hat{a}
=
\arg\max_{a \in \mathcal{A}(\hat{f})}
P(a \mid \mathbf{x},\hat{f}),
\label{eq:architecture}
\end{equation}
where $\mathcal{A}(\hat{f})$ denotes the architectures associated with the
selected family. Finally, variant identification is conditioned on the selected
architecture:
\begin{equation}
\hat{v}
=
\arg\max_{v \in \mathcal{V}(\hat{a})}
P(v \mid \mathbf{x},\hat{a}),
\label{eq:variant}
\end{equation}
where $\mathcal{V}(\hat{a})$ represents the candidate variants belonging to the
selected architecture. This hierarchical formulation preserves the structural
relationships among AI models and enables fingerprinting at progressively finer
levels of model identity. To isolate per-level discriminability, each level is
also evaluated independently using the ground-truth parent.

\subsection{Cross-Device Fingerprinting}
\label{subsec:crossdevice}
In federated learning, the participating clients train the \emph{same} model on
heterogeneous hardware and operating systems. Classifying each client in isolation discards the most
discriminative cue available to a passive observer: how the timing and transfer
shape of the identical model differ \emph{across} devices, which is governed by
the architecture's per-layer compute profile. We therefore represent each
experiment $e$ by a single feature vector that concatenates the per-device
summaries,
\begin{equation}
\mathbf{x}_e =
\big[\,
\mathbf{f}^{(1)}_{e} \,\Vert\,
\mathbf{f}^{(2)}_{e} \,\Vert\,
\mathbf{f}^{(3)}_{e} \,\Vert\,
\mathbf{r}_e
\,\big],
\label{eq:crossdevice}
\end{equation}
where $\mathbf{f}^{(c)}_{e}$ is the size-normalized feature vector of client $c$
and $\mathbf{r}_e$ collects cross-device compute-gap and throughput ratios.
Because~\eqref{eq:crossdevice} yields few high-dimensional examples, family
identification~\eqref{eq:family} is performed with an $\ell_2$-regularized
multinomial logistic regression whose regularization strength is chosen by
nested cross-validation (Sec.~\ref{subsec:classification}); tree ensembles
overfit this regime.

\subsection{Feature Selection, Classification, and Validation}
\label{subsec:classification}
The extracted traffic fingerprints train supervised classifiers for each
hierarchical identification task. Feature discriminability is examined using
Fisher-score-based selection, which evaluates the separation between classes
relative to their within-class variability, and is performed independently
within each training fold to prevent information leakage from held-out
executions. For the per-client representations we use Fisher-selected random
forests; for the cross-device representation~\eqref{eq:crossdevice} we use
regularized logistic regression. Exact-model identification is performed by
nearest-reference matching on the logarithm of the uploaded payload size.

Generalization is assessed under experiment-disjoint, leave-one-model-out
validation, in which every held-out model is absent from training, so
performance measures recognition of unseen models rather than memorization. All
hyperparameter and feature-selection choices are made within an inner loop over
the training models, so held-out executions never influence model selection.
Performance is reported as balanced accuracy, whose chance level is
$1/|\mathcal{C}|$ for a $|\mathcal{C}|$-class task, independent of class sizes;
reliability is quantified by a group-bootstrap $95\%$ confidence interval over
models and a label-permutation test against chance.

At the variant level, each model is realized by a single execution, and variant
identity is recovered through exact-model identification: every variant produces
a distinct serialized size, enabling nearest-reference matching on the uploaded
payload. Confusion matrices and feature discriminability are examined to
characterize residual ambiguity among AI workloads, particularly where different
architectures exhibit similar network communication patterns.

Robustness to size concealment is assessed by simulating a padding defense in
which every transfer is padded up to a size bucket, and by repeating family
identification using only features independent of transfer size: the compute
gap, its dispersion, and cross-device gap ratios.
